\documentclass{article}
\usepackage[T1]{fontenc}
\usepackage{lmodern}
\usepackage{graphicx}
\usepackage{amsmath,amssymb}
\usepackage[a4paper,margin=2cm]{geometry}
\usepackage[absolute,overlay]{textpos}
\setlength{\TPHorizModule}{1mm}
\setlength{\TPVertModule}{1mm}
\usepackage[indonesian]{babel}
\usepackage{hyperref}
\setlength{\emergencystretch}{2em}
% unit: Halaman 12

\begin{document}

% === Halaman 12 (python/native) ===

12

• Selanjutnya Alice menghitung kunci privat d dengan kekongruenan:


ed  1 (mod (n))   → d =

1+𝑘 ∅(𝑛)

𝑒


 d adalah balikan e dalam modulus (n)

    d dapat dihitung dengan dengan rumus:


    Coba  berbagai nilai-nilai k = 1, 2, 3, …, sampai mendapatkan d bilangan bulat: 


k = 1 → d = (1 + 3220  1)/79 = 3221/79 = 40,77    (bukan bilangan bulat)


k = 2 → d = (1 + 3220  2)/79 = 6441/79 = 81,53    (bukan bilangan bulat)


k = … dst


k = 25 → d = (1 + 3220  25)/79 = 80501/79 = 1019    (bilangan bulat)

   diperoleh nilai d yang bulat adalah 1019. Ini adalah kunci privat. 

• Jadi, kunci privat Alice adalah (d, n) = (1019, 3337). Keep them secret!

𝑑=

1+𝑘(𝑛)

𝑒

 = 

1+3220𝑘

79

\end{document}
